\documentclass[12pt,a4paper,twoside,spanish]{report}
\usepackage{graphicx}
\usepackage{array}
\usepackage{enumerate}
\usepackage{graphicx, subfigure}
\usepackage[spanish,activeacute]{babel}
\usepackage[T1]{fontenc}
\usepackage{lmodern}
\usepackage[utf8]{inputenc}
\usepackage[left=2.5cm,top=3.5cm,right=2.5cm,bottom=2.5cm]{geometry}
\usepackage{fancyhdr}
\usepackage{setspace}
\usepackage{url}
\usepackage{sectsty}\allsectionsfont{\usefont{OT1}{phv}{bc}{n}\selectfont}
\setcounter{secnumdepth}{5}
\usepackage{ccaption}
   \captionnamefont{\bfseries\small}
   \captiontitlefont{\small}
   \captiondelim{\hspace{1em}}
   \hangcaption
\usepackage{multirow}
\usepackage{longtable,ltcaption}
\usepackage{pdfpages} % to import PDF pages

\usepackage{relsize}
\usepackage{framed}
\usepackage{amsmath}
\usepackage{color}
%\usepackage[ruled,linesnumbered,portugues]{algorithm2e}

\usepackage{listings}
\definecolor{backgroundcode}{rgb}{0.95,0.95,0.95}


%referencias de la pagina
\usepackage{hyperref}

\newcommand{\anio}{2015}
\newcommand{\numguia}{8}

%\newcommand{\nombreguia}{Introducción}
%\newcommand{\nombreguia}{STL I}
%\newcommand{\nombreguia}{STL II}
%\newcommand{\nombreguia}{Complejidad}
%\newcommand{\nombreguia}{Backtracking}
%\newcommand{\nombreguia}{Grafos I}
%\newcommand{\nombreguia}{Grafos II}
\newcommand{\nombreguia}{Grafos III}
%\newcommand{\nombreguia}{}


%% propiedades del pdf
\hypersetup{
    pdfstartview={FitH},  
    pdftitle={Gu\'ia pr\'actica \numguia -- \nombreguia -- \anio},
    pdfauthor={David N. Campo - Daniel Ambort},     
    pdfsubject={Gu\'ia pr\'actica \numguia - Programaci\'on Competitiva \anio},
    pdfcreator={David N. Campo},   
    pdfproducer={David N. Campo}, 
    pdfkeywords={algorithms, c++, acm icpc, UVa, SPOJ, programming}, 
    colorlinks,
    citecolor=black,
    filecolor=black,
    linkcolor=black,
    urlcolor=blue,
}

% Estilo de página %%%%%%%%%%%%%%%%%%%%
\fancyhead{}
\fancyhead[LO,RE]{\vfill{\footnotesize Programaci\'on Competitiva $\anio$ - Gu\'ia pr\'actica $\numguia$ -- \nombreguia}}
\fancyhead[RO,LE]{\includegraphics[scale=0.15]{../img/logo_UTN_color.pdf}}

\fancyfoot{}
\fancyfoot[RO]{\hfill\includegraphics[scale=0.9]{../img/footlineR.pdf}\newline\footnotesize Página \textbf{\thepage}}
\fancyfoot[LE]{\includegraphics[scale=0.9]{img/footlineL.pdf}\newline\footnotesize Página \textbf{\thepage}}

\pagestyle{fancy}
\fancypagestyle{plain}{
   \fancyhead{}
	\fancyhead[LO,RE]{\vfill{\footnotesize Programaci\'on Competitiva $\anio$ - Gu\'ia pr\'actica $\numguia$ -- \nombreguia}}
   \fancyhead[RO,LE]{\includegraphics[scale=0.15]{../img/logo_UTN_color.pdf}}
}   

\newpage\setcounter{page}{1}
\pagenumbering{arabic}



\begin{document}


\section*{Gu\'ia Pr\'actica \numguia -- \nombreguia}

\begin{enumerate}[1 -]


	\item Registrarse como usuario en los jueces en línea de SPOJ\footnote{\href{http://www.spoj.com/} {http://www.spoj.com/}} y UVa\footnote{\href{http://uva.onlinejudge.org/} {http://uva.onlinejudge.org/}}.
Con respecto al nombre de usuario, debe utilizar un identificador formado por las iniciales de su/s nombre/s seguidas de su apellido completo, todo en minúsculas y finalmente el año actual - en dígitos. Por ejemplo, si Ud. se llama Juan Darío López, le corresponde el nombre de usuario jdlopez\anio.

	\item Descargar alg\'un compilador y enlazador de C++. En el laboratorio est\'a instalado el IDE ZinjaI\footnote{\href{http://zinjai.sourceforge.net/} {http://zinjai.sourceforge.net/}}.
	
	\item Resolver en lenguage C++ y obtener \textbf{ACCEPTED} en los jueces correspondientes para los siguientes problemas:
	
	\begin{enumerate}[a)]
	
		\item Life, The Universe and Everything - \href{http://www.spoj.com/problems/TEST/} {http://www.spoj.com/problems/TEST/}
		\item He is offside! - \href{http://www.spoj.com/problems/OFFSIDE/}{http://www.spoj.com/problems/OFFSIDE/}
		\item Desde Hasta - \href{http://www.spoj.com/problems/DESDEHAS/}{http://www.spoj.com/problems/DESDEHAS/}
		\item Add and Reversed Sums - \href{http://www.spoj.com/problems/ADDREVUT/}{http://www.spoj.com/problems/ADDREVUT/} 
		\item Blowing Fuses - \href{http://www.spoj.com/problems/BLOWFUSE/}{http://www.spoj.com/problems/BLOWFUSE/} 
		\item Tex Quotes - \href{http://uva.onlinejudge.org/external/2/272.html}{http://uva.onlinejudge.org/external/2/272.html} 
		\item Skew Binary - \href{http://uva.onlinejudge.org/external/5/575.html}{http://uva.onlinejudge.org/external/5/575.html}
		\item Minesweeper - \href{http://uva.onlinejudge.org/external/101/10189.html}{http://uva.onlinejudge.org/external/101/10189.html} 	
		\item LC-Display - \href{http://uva.onlinejudge.org/external/7/706.html}{http://uva.onlinejudge.org/external/7/706.html} 	
		\item The 3n + 1 problem - \href{http://uva.onlinejudge.org/external/1/100.html}{http://uva.onlinejudge.org/external/1/100.html} 	
	\end{enumerate}
	
	\item  Conformar equipo de 3 alumnos (con un nombre de fantasía adecuado) y comunicar dichos datos (nombre equipo - integrantes ) a los docentes de la cátedra quienes les asignarán un nombre para el registro de dicho equipo en Spoj.com y UVa.

\end{enumerate}





\end{document}
